\section{Mixed characteristic}\label{sec:mixed}

We work one irreducible component at a time. Finite integral extensions
make the differential entries deep enough to apply the uniform complex
estimate. An upper bound for the resulting normalized length is obtained
from a prime filtration, and is independent of those extensions after a
separate parameter-power limit. The following two multiplicity estimates
supply that upper bound. Whenever an
ideal is used on a quotient ring, its image in that quotient is understood.

\begin{lemma}\label{lem:mixed-quotient}
Let $(A,\mathfrak a)$ be a Noetherian local ring of positive dimension
$t$, let $H$ be $\mathfrak a$-primary, and let $f\in H$.  If
$\dim A/(f)=t-1$, then
\[
  e(H,A)\leq e(H,A/(f)).
\]
\end{lemma}

\begin{proof}
Multiplication by $f$ induces a map $A/H^{j-1}\longrightarrow A/H^j$
whose cokernel is $A/(H^j,f)$.  Consequently
\[
 \length_A(A/(H^j,f))
 \geq \length_A(A/H^j)-\length_A(A/H^{j-1}).
\]
For large $j$ these are Hilbert--Samuel polynomials.  The polynomial
on the right has leading coefficient $e(H,A)/(t-1)!$, and the one
on the left has leading coefficient $e(H,A/(f))/(t-1)!$.
Comparing them proves the assertion, including the case $t=1$.
\end{proof}

\begin{lemma}\label{lem:mixed-asymptotic}
Let $(B,\mathfrak b)$ be a complete Noetherian local domain of
dimension $h$, and let $1\leq d\leq h$. Suppose that an ideal
$L\subseteq\mathfrak b$ and elements $x_1,\ldots,x_d\in\mathfrak b$ satisfy
$\sqrt{L+(x_1,\ldots,x_d)}=\mathfrak b$.  Put
\[
 H_N=L+(x_1^N,\ldots,x_d^N),\qquad N\geq1.
\]
The following assertions hold.
\begin{enumerate}
\item If $\dim B/L<d$, then $e(H_N,B)=O(N^{d-1})$.
\item If $L=\mathfrak q$ is prime, $\dim B/\mathfrak q=d$, and
$B_{\mathfrak q}$ is regular, then
\[
 e(H_N,B)\leq N^d e((x_1,\ldots,x_d),B/\mathfrak q)
                 +O(N^{d-1}).
\]
\end{enumerate}
All constants concern the fixed data $B,L,x$ and are independent of $N$.
\end{lemma}

\begin{proof}
Choose $a\geq1$ with $\mathfrak b^a\subseteq L+(x_1,\ldots,x_d)$,
and set $r_N=d(N-1)+1$.  The pigeonhole principle for monomials gives
\begin{equation}\label{eq:mixed-power-containment}
 \mathfrak b^{a r_N}
 \subseteq \bigl(L+(x_1,\ldots,x_d)\bigr)^{r_N}
 \subseteq L+(x_1^N,\ldots,x_d^N)=H_N.
\end{equation}
In particular, on any fixed nonzero quotient $A$ of $B$ of dimension
$t\leq d-1$, monotonicity of multiplicity under inclusion of primary
ideals gives
\begin{equation}\label{eq:mixed-small-quotient}
 e(H_N,A)\leq e(\mathfrak b^{a r_N},A)
       =(a r_N)^t e(\mathfrak b,A)=O(N^{d-1}).
\end{equation}
This also applies when $t=0$.

To use this estimate, we will pass from $B$ to fixed quotients of
smaller dimension.  Quotienting by an element of $L$ that lowers the
dimension by one can only increase the multiplicity of $H_N$, by
Lemma~\ref{lem:mixed-quotient}.  In the first case we will reach
dimension $d-1$.  In the second case we will first reach dimension
$d$, choosing the quotient so that its localization at $\mathfrak q$
is the residue field of $B_{\mathfrak q}$.  The component
$B/\mathfrak q$ will therefore have coefficient one in the
associativity formula, while every other component of that quotient
will admit one more dimension-lowering quotient.

The ring $B$ is catenary and equidimensional.  Thus, for every prime
$\mathfrak p\subset B$, one has
$\dim B/\mathfrak p=h-\operatorname{ht}\mathfrak p$.
We will also use the following consequence.  A list of $j$ partial
parameters in $B$ has all its minimal primes of height $j$: its
quotient has dimension $h-j$, which gives the lower bound on these
heights, and the generalized principal ideal theorem gives the
upper bound.  In particular its quotient is equidimensional.

For the first assertion, select $h-d+1$ partial parameters
$v_1,\ldots,v_{h-d+1}$ in $L$.  Here is the required prime avoidance
argument.  After $j<h-d+1$ choices, a minimal prime of the preceding
ideal has height $j$ and quotient dimension $h-j\geq d$.  It cannot
contain $L$, because $\dim B/L<d$.  Hence the next element can be
chosen in $L$ outside all these finitely many primes.  Each choice
lowers the dimension by one.  Applying
Lemma~\ref{lem:mixed-quotient} successively, with $H_N$ fixed, gives
\[
 e(H_N,B)\leq e(H_N,B/(v_1,\ldots,v_{h-d+1})).
\]
The quotient on the right is fixed and has dimension $d-1$, so
\eqref{eq:mixed-small-quotient} proves the first assertion.

For the second assertion put $c=h-d=\operatorname{ht}\mathfrak q$.
We choose $v_1,\ldots,v_c\in\mathfrak q$ that are partial parameters
in $B$ and generate the maximal ideal of the regular local ring
$B_{\mathfrak q}$.  To justify both conditions simultaneously,
suppose that $v_1,\ldots,v_j$ have been chosen, with $j<c$, and that
their images in
\[
 \mathfrak qB_{\mathfrak q}/(\mathfrak qB_{\mathfrak q})^2
\]
are linearly independent over $\kappa(\mathfrak q)$.  Let $E_j$ be
the contraction to $B$ of
\[
 (v_1,\ldots,v_j)B_{\mathfrak q}
                 +(\mathfrak qB_{\mathfrak q})^2.
\]
Since this cotangent space has dimension $c$, the ideal
$\mathfrak q$ is not contained in $E_j$.  Nor is it contained in
any minimal prime of $(v_1,\ldots,v_j)$: those primes have height
$j<c$.  Prime avoidance, with the single possibly nonprime ideal
$E_j$, therefore supplies an element $v_{j+1}$ in $\mathfrak q$
outside their union.  This preserves both conditions.  At the end,
Nakayama's Lemma gives
\[
 (v_1,\ldots,v_c)B_{\mathfrak q}=\mathfrak qB_{\mathfrak q}.
\]
For $c=0$ this is the empty list.

Set $A=B/(v_1,\ldots,v_c)$, which is equidimensional of dimension
$d$.  Repeated application of Lemma~\ref{lem:mixed-quotient}, followed
by associativity of multiplicity, gives
\begin{equation}\label{eq:mixed-cut-associativity}
 e(H_N,B)\leq e(H_N,A)
   =\sum_{\mathfrak p\in\operatorname{Min}(A)}
      \length_{B_{\mathfrak p}}(A_{\mathfrak p})\,
      e(H_N,B/\mathfrak p),
\end{equation}
where primes of $A$ are identified with primes of $B$.
The prime $\mathfrak q$ is one of these minimal primes, and its
coefficient is
\[
 \length_{B_{\mathfrak q}}(A_{\mathfrak q})
 =\length_{B_{\mathfrak q}}
      (B_{\mathfrak q}/\mathfrak qB_{\mathfrak q})=1.
\]
On this component $H_N$ becomes $(x_1^N,\ldots,x_d^N)$.
The parameter-ideal power formula therefore makes its contribution
exactly $N^d e((x_1,\ldots,x_d),B/\mathfrak q)$.

For any other prime $\mathfrak p$ in
\eqref{eq:mixed-cut-associativity}, the two distinct primes
$\mathfrak p$ and $\mathfrak q$ have the same height $c$ and are
incomparable.  Choose $w_{\mathfrak p}\in\mathfrak q\setminus
\mathfrak p$.  Its image is nonzero in the complete local domain
$B/\mathfrak p$ of dimension $d$, and its quotient has dimension
$d-1$.  It belongs to the image of every $H_N$.  Thus
Lemma~\ref{lem:mixed-quotient} and
\eqref{eq:mixed-small-quotient} give
\[
 e(H_N,B/\mathfrak p)
 \leq e(H_N,B/(\mathfrak p,w_{\mathfrak p}))=O(N^{d-1}).
\]
The number of these components and their coefficients in
\eqref{eq:mixed-cut-associativity} are fixed and finite.  Summing
proves the second assertion.
\end{proof}

\begin{theorem}\label{thm:mixed-characteristic}
Assume the notation and conclusions of Proposition~\ref{prop:setup},
and suppose that the source domain $R$ has mixed characteristic.
Then
\[
 e(S)\geq e(T)=e(R).
\]
\end{theorem}

\begin{proof}
Write $\mathfrak t$ for the maximal ideal of $T$, and recall that
$\dim T=h$, $\dim S=d>0$, $S=T/J$, and
\[
 F=P\otimes_T K(x_1,\ldots,x_d;T),
\]
where $P$ is a finite free resolution of $S$, of length at most
$h-d$.  The complex $F$ has length at most $h$, has $F_0=T$, and is
contractible away from the closed point.  Its differentials have
entries in $\mathfrak t$, and its ranks have alternating sum zero.
The images of $x_1,\ldots,x_d$ form a minimal reduction of the maximal
ideal of $S$; in particular $\sqrt{J+(x_1,\ldots,x_d)}=\mathfrak t$.
The residue field, denoted by $k$, is algebraically closed.

Fix a minimal prime $P_0$ of $T$ and put
$B=T/P_0$, with maximal ideal $\mathfrak b$.  By
Proposition~\ref{prop:setup}, $B$ is a complete local domain of mixed
characteristic, has dimension $h$, and has residue field $k$.
Choose a parameter ideal $(u_1,\ldots,u_h)$ that reduces
$\mathfrak b$.

\medskip
\noindent\emph{Finite extensions that make the differentials deep.}
Fix an integer $q>h$ and choose roots $z_i=u_i^{1/q}$ in an algebraic
closure of $\operatorname{Frac}(B)$.  The ring
$B'=B[z_1,\ldots,z_h]$ is a finite domain extension of $B$.
It is complete local, has residue field $k$, and has dimension $h$.
The ideal $(z_1,\ldots,z_h)$ is primary for its maximal ideal.
Since $\mathfrak b$ is integral over $(u_1,\ldots,u_h)$, every
element of $\mathfrak b$ is integral over $(z_1,\ldots,z_h)^qB'$.

Choose $a\geq1$ with $p^a\in(z_1,\ldots,z_h)B'$, where $p$ is
the residue characteristic.  Apply Heitmann's integral-extension
form of the Brian\c{c}on--Skoda Theorem~\cite{Heitmann}, as stated in
\cite[Proposition~2.8]{HeitmannMa} to the generating list
\[
 p^a,z_1,\ldots,z_h
\]
of the ideal $I=(z_1,\ldots,z_h)B'$.  This list has $n=h+1$
members and satisfies the theorem's hypothesis
$p\in\sqrt{(p^a,z_1)B'}$.  With $k'=q-h-1\geq0$, its conclusion
sends each element integral over $I^{n+k'}=I^q$ into
$I^{k'+1}=I^{q-h}$ after an integral extension.
Here redundant generators are allowed by the theorem's statement.
Apply it to a finite generating set of $\mathfrak b$.  Taking a
prime above zero in each extension if necessary gives integral
domain extensions, which can be embedded in a common algebraic
closure of $\operatorname{Frac}(B)$.  Adjoin the finitely many
integral coefficients witnessing the resulting memberships.
We obtain a domain $D$ finite over $B$ such that
\begin{equation}\label{eq:mixed-deep-entries}
 \mathfrak bD\subseteq(z_1,\ldots,z_h)^{q-h}D.
\end{equation}
We may replace $D$ by its normalization: it remains finite by
excellence, and the inclusion persists.  A finite algebra over a
complete local ring is a product of complete local rings; since
$D$ is a domain, there is just one factor.  Thus $D$ is complete
local of dimension $h$, and its residue field is $k$ because it
is finite over the algebraically closed field $k$.

Put $r=\rank_B D$.  For every $\mathfrak b$-primary ideal $H$,
the finite-module multiplicity formula and equality of residue
fields give
\begin{equation}\label{eq:mixed-rank-formula}
 e(HD,D)=r e(H,B).
\end{equation}
Indeed, a $\operatorname{Frac}(B)$-basis of
$D\otimes_B\operatorname{Frac}(B)$ can be
cleared of denominators to give an injection $B^r\longrightarrow D$
with torsion cokernel.  This cokernel has dimension less than $h$,
so it does not contribute to multiplicity.  Moreover, lengths of
finite-length $D$-modules are the same over $B$ and over $D$ because
the residue fields agree.  Applying
\eqref{eq:mixed-rank-formula} to $(u_1,\ldots,u_h)$ and using the
parameter-ideal power formula gives
\begin{equation}\label{eq:mixed-mu-normalization}
 q^h e((z_1,\ldots,z_h),D)
 =e((u_1,\ldots,u_h),D)=r e(B).
\end{equation}

\medskip
\noindent\emph{A lower bound for the normalized length.}
Use the algebra $C$ and additive length $\lambda$ supplied for $D$
by Theorem~\ref{thm:length-inputs}, and put
$\mu=e((z_1,\ldots,z_h),D)$.  Extend $F$ along
$T\longrightarrow B\longrightarrow D\longrightarrow C$.
Its entries belong to $(z_1,\ldots,z_h)^{q-h}C$ by
\eqref{eq:mixed-deep-entries}.  Local contractions survive this
base change.  The ideals $\mathfrak tD$ and $(z_1,\ldots,z_h)D$
have the same radical, so the required support is also preserved.
Theorem~\ref{thm:complex}, with $s=q-h$ and zeroth rank one,
therefore gives a constant $K$ depending only on $h$ and the fixed
ranks of $F$ such that
\[
 \lambda(C/(J,x_1,\ldots,x_d)C)
 \geq \mu(q-h)^h\left(1-\frac{K}{q-h}\right)
\]
for all sufficiently large $q$.  Together with
\eqref{eq:mixed-mu-normalization}, this reads
\begin{equation}\label{eq:mixed-lower-bound}
 \frac{\lambda(C/(J,x_1,\ldots,x_d)C)}{r}
 \geq e(B)\left(1-\frac hq\right)^h
                  \left(1-\frac{K}{q-h}\right).
\end{equation}
In particular, $K$ is independent of the choices of $D$, $C$, and
$\lambda$ as $q$ varies.

\medskip
\noindent\emph{The power rule for the quotient lengths.}
For now keep $q$, $D$, $C$, and $\lambda$ fixed.  For every list of
positive integers $a_1,\ldots,a_d$, consider
\[
 F(a)=P_C\otimes_C K(x_1^{a_1},\ldots,x_d^{a_d};C).
\]
Before base change this complex is contractible away from the
closed point: outside $V(J)$ the resolution $P$ is split exact,
and outside $V(x_1,\ldots,x_d)$ the Koszul complex is contractible.
Its length is at most $h$.  Lemma~\ref{lem:acyclicity} therefore
shows that its positive homology has $\lambda$-length zero and
that its zeroth homology has finite $\lambda$-length.  Thus
\[
 \chi_\lambda(F(a))
 :=\sum_i(-1)^i\lambda(H_i(F(a)))
 =\lambda(C/(J,x_1^{a_1},\ldots,x_d^{a_d})C).
\]

To check the power rule without any finiteness assumption on
partial Koszul complexes, fix all exponents but $a_i$, and let
$E$ denote $P_C$ tensored with the remaining Koszul factors.
For positive integers $v,w$, the composable scalar maps
$x_i^v,x_i^w$ on $E$ give the mapping-cone triangle
\[
 \operatorname{Cone}(x_i^v:E\to E)
 \longrightarrow\operatorname{Cone}(x_i^{v+w}:E\to E)
 \longrightarrow\operatorname{Cone}(x_i^w:E\to E)
 \longrightarrow\operatorname{Cone}(x_i^v:E\to E)[1].
\]
Every cone includes a positive power of each $x_j$, so all three
have finite homology lengths by the preceding argument.  Additivity
of $\lambda$ in their long exact homology sequence gives additivity
of Euler characteristic.  Hence that Euler characteristic is
linear in each positive exponent.  Applying this to all $d$
exponents proves
\begin{equation}\label{eq:mixed-length-power-rule}
 \lambda(C/(J,x_1^N,\ldots,x_d^N)C)
   =N^d\lambda(C/(J,x_1,\ldots,x_d)C)
 \qquad(N\geq1).
\end{equation}
No assertion that the homology of $E$ itself has finite length was
needed.

\medskip
\noindent\emph{An upper bound from a prime filtration.}
Fix a finite prime filtration
\[
 0=S_0\subset S_1\subset\cdots\subset S_m=S,
 \qquad S_j/S_{j-1}\simeq T/Q_j.
\]
All $Q_j$ contain $J$; the list records repetitions.  Set
\[
 M_N=C/(x_1^N,\ldots,x_d^N)C,
 \qquad H_{j,N}=Q_jB+(x_1^N,\ldots,x_d^N)B.
\]
Tensoring each filtration sequence with $M_N$ gives a right exact
sequence
\[
 S_{j-1}\otimes_T M_N\longrightarrow S_j\otimes_T M_N
                   \longrightarrow C/H_{j,N}C\longrightarrow0.
\]
Every module here is killed by $J+(x_1^N,\ldots,x_d^N)$, whose
extension to $D$ is primary for its maximal ideal.  Their lengths
are finite: for example, each $S_j\otimes_T M_N$ is a quotient of
a finite direct sum of $C/(J,x_1^N,\ldots,x_d^N)C$.
The image of the first arrow is a quotient of its source, so
additivity and nonnegativity of $\lambda$ yield
\[
 \lambda(S_j\otimes_T M_N)
 \leq\lambda(S_{j-1}\otimes_T M_N)+\lambda(C/H_{j,N}C).
\]
Summing these inequalities, applying the primary-ideal bound from
Theorem~\ref{thm:length-inputs}, and then applying
\eqref{eq:mixed-rank-formula}, we obtain
\begin{equation}\label{eq:mixed-filtration-bound}
 \lambda(C/(J,x_1^N,\ldots,x_d^N)C)
 \leq\sum_{j=1}^m\lambda(C/H_{j,N}C)
 \leq r\sum_{j=1}^m e(H_{j,N},B).
\end{equation}
Thus this use of the prime filtration requires only right
exactness of tensor product, not flatness of $C$ over $T$.

We apply Lemma~\ref{lem:mixed-asymptotic} to each term of the last
sum.  Since $Q_j\supseteq J$, one has $\dim T/Q_j\leq d$.
If $\dim T/Q_j<d$, then also $\dim B/Q_jB<d$.  If $P_0$ is not
contained in $Q_j$, its image in the domain $T/Q_j$ is nonzero,
and hence
\[
 \dim B/Q_jB=\dim T/(P_0+Q_j)<\dim T/Q_j\leq d.
\]
In either case the first assertion of the Lemma gives
$e(H_{j,N},B)=O(N^{d-1})$.

The remaining case is $\dim T/Q_j=d$ and $P_0\subseteq Q_j$.
Such a $Q_j$ is minimal over $J$: a strictly smaller prime over
$J$ would give a quotient of dimension greater than $d$.
The generic regularity in Proposition~\ref{prop:setup} gives that
$T_{Q_j}$ is regular.  It follows that $Q_j$ contains exactly one
minimal prime of $T$, namely $P_0$, and
\[
 B_{Q_jB}=T_{Q_j}.
\]
Indeed the localization of $P_0$ is the unique minimal prime of
the regular local domain $T_{Q_j}$ and is therefore zero.
The second assertion of Lemma~\ref{lem:mixed-asymptotic} now gives
\begin{equation}\label{eq:mixed-component-asymptotic}
 e(H_{j,N},B)
 \leq N^d e((x_1,\ldots,x_d),T/Q_j)+O(N^{d-1}).
\end{equation}

\medskip
\noindent\emph{The two limits and the component sum.}
Divide \eqref{eq:mixed-filtration-bound} by $rN^d$ and use
\eqref{eq:mixed-length-power-rule}.  Let $N$ tend to infinity
while keeping the chosen $q,D,C,\lambda$ fixed.  The preceding
asymptotic estimates imply
\begin{equation}\label{eq:mixed-uniform-upper-bound}
 \frac{\lambda(C/(J,x_1,\ldots,x_d)C)}{r}
 \leq
 \sum_{\substack{1\leq j\leq m\,:\ \dim T/Q_j=d\\P_0\subseteq Q_j}}
       e((x_1,\ldots,x_d),T/Q_j).
\end{equation}
The right side does not involve any auxiliary extension or length
algebra.  Combine this bound with
\eqref{eq:mixed-lower-bound} and then let $q$ tend to infinity.
Uniformity of $K$ gives
\begin{equation}\label{eq:mixed-one-component}
 e(T/P_0)\leq
 \sum_{\substack{1\leq j\leq m\,:\ \dim T/Q_j=d\\P_0\subseteq Q_j}}
       e((x_1,\ldots,x_d),T/Q_j).
\end{equation}
There is no exchange of limits: the $N$-limit was taken separately
for each fixed choice of the data belonging to $q$.

Finally sum \eqref{eq:mixed-one-component} over the minimal primes
$P_0$ of $T$.  The ring $T$ is equidimensional and generically
reduced by Proposition~\ref{prop:setup}; its generic local lengths
are therefore all one, and associativity gives
\[
 e(T)=\sum_{P_0\in\operatorname{Min}(T)}e(T/P_0).
\]
Every $Q_j$ of dimension $d$ belongs to exactly one of these
components, because $T_{Q_j}$ is regular.  Thus summation counts
each occurrence of $Q_j$ in the filtration exactly once.  The
number of occurrences of a fixed such prime $Q$ is
$\length_{T_Q}(S_Q)$, as seen by localizing the filtration at $Q$.
The multiplicity formula for $S$ consequently yields
\[
 e(T)
 \leq\sum_{\substack{1\leq j\leq m\,:\ \dim T/Q_j=d}}
           e((x_1,\ldots,x_d),T/Q_j)
 =e((x_1,\ldots,x_d),S)=e(S).
\]
The equality $e(T)=e(R)$ from Proposition~\ref{prop:setup}
finishes the proof.
\end{proof}
